El sistema de navegació europeu Galileu estarà format per trenta satèl·lits distribuïts en tres plans orbitals a $2,36\times 10^{4}\ \mathrm{km}$ d'altura sobre la Terra, i cada un d'ells descriurà una òrbita circular. Calculeu:

\figura[0.45]{fig1.png}
{\centering\textit{Esquema de les òrbites dels satèl·lits Galileu}\par}

\begin{apartat}{1}
Quin període de rotació tindran aquests satèl·lits?
\end{apartat}

\begin{apartat}{1}
Quina serà la velocitat orbital dels satèl·lits?
\end{apartat}

\dades{%
  $G = 6,67\times 10^{-11}\ \mathrm{N\,m^2\,kg^{-2}}$\\
  $R_\mathrm{Terra} = 6,37\times 10^{6}\ \mathrm{m}$\\
  $M_\mathrm{Terra} = 5,98\times 10^{24}\ \mathrm{kg}$}
